\section{評価実験}

提案手法を二つの観点から評価する．
一つは，学習済みモデル単体の性能を評価セット上で測る\textbf{オフライン評価}であり，
候補の正解一致率（Acc@$k$）と入力効率（KSPC）を指標に，
文脈・ビーム幅・LoRA の寄与をアブレーションで切り分ける．
もう一つは，実利用者の操作を要する\textbf{実ユーザ実験}であり，
母音入力に特有の認知負荷を計測 Web アプリを用いて測定する．
本章では評価指標（\ref{subsec:metrics}節）と比較対象（\ref{subsec:compare}節），
実ユーザ実験の設計（\ref{subsec:user}節）を述べ，結果は次章にまとめる．

\subsection{評価指標}
\label{subsec:metrics}
候補提示型の入力支援であることを踏まえ，Top-1 の完全一致だけでなく
Top-3 / Top-5 を重視する．母音入力 $\bm{v}_i$，正解 $y_i$，上位 $k$ 候補集合 $\mathcal{Y}_k(\bm{v}_i)$ に対し，
\begin{equation}
  \mathrm{Acc@}k = \frac{1}{N}\sum_{i=1}^{N}\mathbf{1}\!\left[y_i \in \mathcal{Y}_k(\bm{v}_i)\right]
\end{equation}
を用いる．

入力効率は KSPC（Keystrokes Per Character）で評価する．
\begin{equation}
  \mathrm{KSPC} = \frac{K}{C}
\end{equation}
ここで $K$ は打鍵数，$C$ は正解文字数である．
本稿では誤りなしで入力し切るのに必要な最小打鍵数を $K$ とする理想 KSPC を用い，
提案手法（$K$＝母音キー数）と既存IME方式（ローマ字＝訓令式，JIS かな，フリック，トグル）を
同一の正解文字数 $C$ のもとで比較する．

\subsection{比較対象と実装状況}
\label{subsec:compare}
入力効率（KSPC）については，提案手法（母音入力）に加え，
既存IME方式（ローマ字＝訓令式，JIS かな入力，フリック入力，トグル入力）の
理想 KSPC（誤りなしの最小打鍵数 $\div$ 正解文字数）を
同一評価セット上から自動算出する．
各方式の打鍵数は正解文の読み（pykakasi により取得）から数え，
分母はどの方式も素の正解文字数 $\mathrm{len}(\text{gold})$ で共通とする．
数字・英字を含む文は母音化で情報が脱落し KSPC が不当に小さくなるため評価対象から除外する．
算出結果は次章の表\ref{tab:input}に示す．

\subsection{実ユーザ実験（認知負荷評価）}
\label{subsec:user}
実利用者の操作を要する評価として，参加者 5 名を対象に実験を行った。
本実験用に，スマートフォン／PC のブラウザから実行する
計測 Web アプリを実装した。
各参加者は 4 つの入力方式（ひらがな入力，母音入力，
標準IME，母音入力＋LLM 推論）で課題文（計 9 文）を入力した。

\subsubsection{評価の目的}
本評価の目的は，母音入力に特有の認知負荷を 2 種類に分離して測定することである。
提案手法である母音入力方式では，従来の日本語 IME には存在しない認知的処理が要求される。
具体的には，利用者は入力したい文章を母音列へ変換し，その母音列を保持しながら入力を行う必要がある。
そのため本研究では，入力精度や入力速度だけでなく，利用時の認知負荷についても評価を行う。
認知負荷の評価にあたっては，広く利用されている主観評価手法である
NASA-TLXを参考にした。
ただし本研究では母音入力特有の負荷を分析するため，一般的な精神的負荷を評価するだけでなく，
\textbf{母音変換負荷}（提示文を母音列へ変換する音韻変換の負荷）と
\textbf{記憶保持負荷}（母音列を保持しながら入力するワーキングメモリ負荷）を
独立した評価項目として設計した点に特徴がある。
この分離により，母音入力の負担がどこから生じるのかを切り分けて議論できる。

\subsubsection{リサーチクエスチョン}
\begin{itemize}
  \item \textbf{RQ1}：母音入力は通常 IME と比較して入力効率を向上させるか。
  \item \textbf{RQ2}：母音入力における認知負荷はどこから生じるか（母音変換負荷か記憶保持負荷か）。
  \item \textbf{RQ3}：文章長の増加によって，母音変換負荷と記憶保持負荷はどのように変化するか。
\end{itemize}

\subsubsection{実験設計}
実験は次の 2 種類を実施する。
\paragraph{(1) 従来 IME との比較（RQ1）}
提案手法の有効性を検証するため，通常の日本語 IME 入力と母音入力を比較する。
評価項目は入力時間，正答率，および主観評価とする。

\paragraph{(2) 母音入力内での認知負荷分析（RQ2，RQ3）}
母音入力方式において，文章長による認知負荷の変化を分析する。
課題文を Level~1：単語（例「おはよう」「電話」），
Level~2：10 文字以内の短文（例「今から帰る」「気をつけてね」），
  Level~3：一般文（例「来週はパソコンを買いに行く」）
の 3 段階に分類した。
本設計の要点は，記憶保持負荷を意図的に課すために，
課題文を一定時間提示した後に画面から隠し，記憶した状態で入力させる点にある。
これにより，文を見ながら入力する場合には現れない記憶保持の負荷を顕在化させ，
母音変換負荷と分離して観察できる。

\subsubsection{客観指標と負荷の対応}
計測アプリは試行ごとに次を記録する。これらは 2 種類の認知負荷の代理指標を与える。
\begin{itemize}
  \item 初動時間 $T_{\mathrm{init}}$（課題文を隠した後から最初の打鍵まで）と
        母音化誤り率 $\mathrm{CER}_{\mathrm{vowel}}$
        （被験者が打った母音列と正解母音列\footnote{母音化規則により算出。}の
        編集距離を参照長で正規化した値）は
        \textbf{母音変換負荷}を反映する。
  \item 総入力時間 $T_{\mathrm{total}}$，正答率，候補選択順位は総合的な効率を反映する。
\end{itemize}

\subsubsection{主観評価項目}
各タスク終了後に，以下の項目について 7 段階評価
（1：全くそう思わない～7：非常にそう思う）を実施する。
母音入力条件（7 項目）：(1) 提示された文章を母音列に変換するのは難しかった（母音変換負荷），(2) 母音列を覚えながら入力するのは難しかった（記憶保持負荷），(3) 入力操作そのものは難しかった，(4) 入力中に混乱やストレスを感じた，(5) 入力後に疲れを感じた，(6) 通常の日本語入力より効率的だと感じた，(7) 日常的に利用したいと思った。
通常 IME 条件（3 項目）：(1) 入力操作そのものは難しかった，(2) 入力中に混乱やストレスを感じた，(3) 入力後に疲れを感じた。
このように，母音変換負荷と記憶保持負荷を独立した項目として測ることで，
単一の精神的負荷尺度では区別できない負荷の由来を分離して分析できる。

\subsubsection{解析方針}
参加者数は 5 名と限られるため，統計的検定によって母集団レベルの有意差を
主張することは本研究の射程としない。代わりに，各方式・各 Level の平均値と
参加者ごとの傾向にもとづいて記述的に分析し，
母音入力に固有の負荷構造を観察することを目的とする。
特に，総入力時間 $T_{\mathrm{total}}$・初動時間 $T_{\mathrm{init}}$・打鍵数と，
母音変換負荷・記憶保持負荷の主観評価について，
推論なし条件（ひらがな入力 対 母音入力）と
推論あり条件（標準 IME 対 母音入力＋推論）に分けて比較し，
母音入力のボトルネックが音韻変換と記憶保持のいずれにあるかを議論する。
